% =============================================================================
% core/page.tex — geometria da página
% -----------------------------------------------------------------------------
% Medidas extraídas do PDF exportado pelo Google Docs (1 bp = 1 pt do Docs).
% =============================================================================

\RequirePackage[
  paperwidth  = 596bp,
  paperheight = 842bp,
  left        = 56.69bp,   % 2 cm
  top         = 28.35bp,   % 1 cm
  textwidth   = 481.9bp,
  textheight  = 786.55bp,  % margem inferior de 27.1bp
  headheight  = 0pt,
  headsep     = 0pt,
  footskip    = 0pt,
  nomarginpar,
]{geometry}

\pagestyle{empty}
\raggedbottom

% Nada de espaçamento automático do TeX: todo espaço vertical é controlado
% pelo modelo de parágrafo em core/paragraph.tex.
\setlength\parindent{0pt}
\setlength\parskip{0pt}
\setlength\topskip{0pt}
\setlength\baselineskip{0pt}
\setlength\lineskip{0pt}
\setlength\lineskiplimit{0pt}
\setlength\normallineskip{0pt}
\setlength\normallineskiplimit{0pt}
\setlength\maxdepth{0pt}
\makeatletter
\setlength\@maxdepth{0pt}
\makeatother

% Sem hifenização (o Docs não hifeniza) e espaçamento uniforme.
\frenchspacing
\hyphenpenalty    = 10000
\exhyphenpenalty  = 10000
\hbadness         = 10000
\hfuzz            = 1bp

% Linhas órfãs/viúvas = 2 (padrão do Google Docs).
\clubpenalty      = 10000
\widowpenalty     = 10000
\interlinepenalty = 0
